\documentclass[11pt]{article}
\usepackage{hyperref}
\usepackage{enumitem}
\usepackage{cmap}
\usepackage[utf8]{inputenc}
\usepackage[T1]{fontenc}
\usepackage[english]{babel}
\usepackage[left=1.2cm,top=1.4cm,right=1.2cm,bottom=1.2cm]{geometry}
\setlength{\parindent}{0pt}
\setlist[itemize]{leftmargin=*, noitemsep, topsep=2pt, partopsep=0pt, parsep=2pt, label={\textbf{-}}}
\begin{document}
\begin{center}
    {\LARGE \textbf{Kevin Bravo}}\\
    Software Engineer | Sistemas de IA, Backend \& Full-Stack \\
    \hrulefill
\end{center}
\begin{center}
    \href{mailto:0bravokevin@gmail.com}{0bravokevin@gmail.com}
    \textbullet\ 
    \href{http://www.kevinbravo.com}{www.kevinbravo.com}
\end{center}
\vspace{4pt}
\begin{center}
    \textbf{Perfil}
\end{center}
\vspace{-4pt}
Ingeniero de software con 4+ años de experiencia construyendo plataformas productivas, sistemas internos y productos full-stack en startups, proyectos públicos y organizaciones sin fines de lucro. He trabajado en arquitectura backend, integración de servicios, modelado de datos, automatización de flujos y ejecución end-to-end de productos. Hoy estoy especialmente enfocado en IA aplicada, sistemas agentic, tooling para desarrolladores y automatizaciones que mejoran operaciones reales.
\vspace{8pt}
\begin{center}
    \textbf{Experiencia}
\end{center}
\vspace{-4pt}
\textbf{Proyecto Educa} \hfill \textit{Remote, Venezuela} \\
\textbf{Chief Technology Officer} \hfill \textit{Dec 2025 -- Present}
\begin{itemize}
    \item Defino la dirección técnica y de producto para iniciativas digitales orientadas a educación, traduciendo necesidades ambiguas en arquitectura, prioridades y planes de ejecución.
    \item Lidero la ejecución hands-on de plataforma, flujos de desarrollo y decisiones de sistema, con responsabilidad directa sobre lo que se construye y cómo evoluciona.
\end{itemize}
\vspace{10pt}
\textbf{Free2Z} \hfill \textit{Remote, Venezuela} \\
\textbf{Principal Software Engineer} \hfill \textit{Apr 2025 -- Apr 2026}
\begin{itemize}
    \item Dirigí la migración técnica de una plataforma productiva de React a SvelteKit, rediseñando la arquitectura central para mejorar mantenibilidad, performance y manejo de estado.
    \item Integré una API en Python/Django con el nuevo frontend, incorporé video en vivo con Dyte, construí un CMS propio y mejoré SSR para fortalecer el desempeño en producción.
\end{itemize}
\vspace{10pt}
\textbf{U.S. Department of State (via AVAA)} \hfill \textit{Remote, Venezuela} \\
\textbf{Software Engineer | Full-Stack Developer} \hfill \textit{Nov 2024 -- Aug 2025}
\begin{itemize}
    \item Construí Billingua Talent como único ingeniero, cubriendo el ciclo completo desde arquitectura de datos y flujos de producto hasta entrega en producción para la Venezuela Affairs Unit del U.S. Department of State.
    \item Diseñé perfiles de candidatos, dashboards internos, flujos de gestión de servicios y vistas de reporting que mejoraron la visibilidad del talento y la toma de decisiones operativas.
\end{itemize}
\vspace{10pt}
\textbf{Freelance} \hfill \textit{Remote} \\
\textbf{Software Engineer} \hfill \textit{Sep 2024 -- Dec 2024}
\begin{itemize}
    \item Entregué plataformas web, scripts de automatización y soluciones end-to-end para clientes de distintos sectores, moviéndome con rapidez desde requerimientos hasta producto usable.
    \item Llevé proyectos desde discovery y priorización hasta desarrollo, lanzamiento, iteración y soporte técnico, muchas veces definiendo la estructura antes de empezar a construir.
\end{itemize}
\vspace{10pt}
\textbf{Asociacion Venezolano Americana de Amistad (AVAA)} \hfill \textit{Remote, Venezuela} \\
\textbf{Software Developer} \hfill \textit{Aug 2023 -- Sep 2024}
\begin{itemize}
    \item Construí SEP, la fuente de verdad para registros de becarios, actividades, horas de voluntariado, asistencia y reportes en 3 capítulos de AVAA, reemplazando flujos fragmentados en hojas de cálculo por un sistema operativo en vivo.
    \item Desarrollé SEA/AES para EducationUSA, centralizando historiales de estudiantes, actividades de outreach, servicios y alianzas institucionales en una plataforma operativa backend-driven.
\end{itemize}
\vspace{10pt}
\textbf{Asociacion Venezolano Americana de Amistad (AVAA)} \hfill \textit{Caracas, Venezuela} \\
\textbf{Software Developer Intern} \hfill \textit{Oct 2022 -- Aug 2023}
\begin{itemize}
    \item Construí un dashboard interactivo en Power BI que convirtió registros dispersos de alumni en un flujo analítico más claro para AVAA.
    \item Escribí scripts en Python para extraer, normalizar y preparar datos históricos, reduciendo limpieza manual y mejorando la confiabilidad del reporting.
\end{itemize}
\vspace{8pt}
\begin{center}
    \textbf{Educación}
\end{center}
\vspace{-4pt}
\textbf{Instituto de Estudios Superiores de Administracion (IESA)} \hfill \textit{Caracas, Venezuela} \\
Maestría en Administración Pública \hfill \textit{En progreso}
\vspace{8pt}
\begin{center}
    \textbf{Liderazgo \& Voluntariado}
\end{center}
\vspace{-4pt}
\textbf{JA Venezuela Alumni} \hfill \textit{Caracas, Venezuela} \\
\textbf{Vice President} \hfill \textit{Mar 2026 -- Present}
\begin{itemize}
    \item Apoyo la planificación estratégica, alianzas y crecimiento institucional de la red nacional de alumni.
    \item Coordino iniciativas enfocadas en emprendimiento, liderazgo profesional y engagement de comunidad.
\end{itemize}
\vspace{10pt}
\textbf{Asociacion Venezolano Americana de Amistad (AVAA)} \hfill \textit{Remote, Venezuela} \\
\textbf{Education Committee Member} \hfill \textit{Sep 2024 -- Present}
\begin{itemize}
    \item Soy administrador principal de PES y doy soporte a operaciones de seguimiento del programa en 3 capítulos de AVAA.
    \item Contribuyo a rediseño curricular, estrategia de fundraising y planificación de largo plazo para Programa Excelencia.
\end{itemize}
\vspace{10pt}
\textbf{Technovation Girls Venezuela} \hfill \textit{Caracas, Venezuela} \\
\textbf{Code Mentor} \hfill \textit{Jan 2026 -- May 2026}
\begin{itemize}
    \item Mentoreo a 4 estudiantes en diseño y desarrollo de una aplicación móvil con MIT App Inventor.
    \item Guío al equipo en investigación, prototipado, resolución técnica y preparación del pitch.
\end{itemize}
\vspace{10pt}
\textbf{Dev3pack} \hfill \textit{Caracas, Venezuela} \\
\textbf{Hackathon Hub Organizer} \hfill \textit{Apr 2026 -- May 2026}
\begin{itemize}
    \item Organicé el hub de Caracas de un hackathon global de Web3 e IA, coordinando logística, alianzas, speakers, mentores y soporte a participantes.
    \item Combiné operaciones, comunidad y mentoría técnica para crear una experiencia útil para builders y estudiantes locales.
\end{itemize}
\vspace{8pt}
\begin{center}
    \textbf{Habilidades}
\end{center}
\vspace{-4pt}
\textbf{Técnicas:} TypeScript, JavaScript, Python, Django, React, SvelteKit, Next.js, SQL, Docker, PostgreSQL, Tailwind CSS, Power BI, ingeniería backend, sistemas de IA, automatización de flujos, APIs e integraciones, modelado de datos, tooling para agentes, scripting y arquitectura de software.
\\[4pt]
\textbf{Profesionales:} Ejecución con alto ownership, liderazgo técnico, product ownership, diseño de sistemas, colaboración cross-functional, mentoring, mejora de procesos, trabajo efectivo en ambigüedad, y construcción de cero a producción. \\[4pt]
\textbf{Idiomas:} Español (Nativo), Inglés (C1). \\[4pt]
\textbf{Intereses:} IA aplicada, sistemas agentic, tooling para desarrolladores, automatización, sistemas distribuidos, infraestructura de software y educación.
\end{document}
